% RxJS Operator Behavior via Cω Type Notation and Mealy-Machine 6-Tuples
% Compile with:  pdflatex operators.tex   (twice, for cross-references)
\documentclass[11pt]{article}

\usepackage[margin=1in]{geometry}
\usepackage{amsmath,amssymb,amsthm}
\usepackage{mathtools}
\usepackage{array}
\usepackage{booktabs}
\usepackage{xcolor}
\usepackage[hidelinks]{hyperref}

% ---- Notation macros ----
\newcommand{\nx}[1]{\mathtt{n}(#1)}      % next
\newcommand{\cmpl}{\mathtt{c}}           % complete
\newcommand{\err}{\mathtt{e}}            % error
\newcommand{\tick}{\tau}                 % clock tick
\newcommand{\emptyout}{\langle\,\rangle} % empty output
\newcommand{\outseq}[1]{\langle #1 \rangle}
\newcommand{\Sstate}{S}
\newcommand{\Sinit}{S_0}
\newcommand{\inalph}{\Sigma}
\newcommand{\outalph}{\Lambda}
\newcommand{\trans}{T}
\newcommand{\outfun}{G}
\newcommand{\card}[1]{\mathsf{#1}}       % cardinality markers

\theoremstyle{definition}
\newtheorem{definition}{Definition}
\newtheorem{theorem}{Theorem}

\title{RxJS Operator Behavior\\
       \large via Erik Meijer's C$\omega$ Stream-Type Notation\\
       and the Mealy-Machine 6-Tuple}
\author{}
\date{}

\begin{document}
\maketitle

\begin{abstract}
We give a dual-lens formalization of RxJS operators. The \emph{static} lens is Erik Meijer's
C$\omega$ stream-type notation, whose cardinality markers $T!$, $T?$, $T^{+}$, $T^{*}$ describe
the shape of a stream. The \emph{dynamic} lens is the Mealy-machine 6-tuple
$M=(S,S_0,\Sigma,\Lambda,T,G)$, generalized so that the output function
$G:S\times\Sigma\to\Lambda^{*}$ emits a \emph{sequence} of symbols. We show that the C$\omega$
cardinality marker is the static projection of the length behavior of $G$.
\end{abstract}

\section{Stream types (C$\omega$)}
\begin{center}
\begin{tabular}{llll}
\toprule
Marker & Name & Interval & RxJS analog \\
\midrule
$T!$        & exactly one  & $[1,1]$      & one \texttt{next} then \texttt{complete} \\
$T?$        & zero or one  & $[0,1]$      & maybe-empty single value \\
$T^{+}$     & one or more  & $[1,\infty)$ & guaranteed non-empty stream \\
$T^{*}$     & zero or more & $[0,\infty)$ & general \texttt{Observable<T>} \\
\bottomrule
\end{tabular}
\end{center}

\section{The Mealy machine}
\begin{definition}[Generalized Mealy machine]
A machine is a 6-tuple $M=(\Sstate,\Sinit,\inalph,\outalph,\trans,\outfun)$ with state space
$\Sstate$, initial state $\Sinit\in\Sstate$, input alphabet $\inalph$, output alphabet
$\outalph$, transition function $\trans:\Sstate\times\inalph\to\Sstate$, and output function
\[
  \outfun : \Sstate\times\inalph \;\longrightarrow\; \outalph^{*}.
\]
The use of $\outalph^{*}$ (rather than $\outalph$) permits zero or more output symbols per
input symbol.
\end{definition}

The shared input alphabet encodes the Rx grammar $(\texttt{next})^{*}(\texttt{error}\mid
\texttt{complete})^{?}$:
\[
  \inalph = \{\,\nx{v} : v\in T\,\}\ \cup\ \{\,\cmpl,\ \err\,\}.
\]
Time-based operators add a clock symbol: $\inalph = \{\nx{v},\cmpl,\err\}\cup\{\tick\}$.

\section{Core operators}

\subsection*{\texttt{map} : $(T\to U)\to T^{*}\to U^{*}$}
\[
\begin{aligned}
\outfun(\text{Active},\nx{v}) &= \outseq{\nx{f(v)}}, &
\outfun(\text{Active},\cmpl)  &= \outseq{\cmpl}, &
\outfun(\text{Active},\err)   &= \outseq{\err}.
\end{aligned}
\]
$|\outfun|\equiv 1$: cardinality preserved.

\subsection*{\texttt{filter} : $(T\to\mathbb{B})\to T^{*}\to T^{*}$,\quad $T!\to T?$}
\[
\outfun(\text{Active},\nx{v}) =
\begin{cases}
\outseq{\nx{v}} & p(v)=\text{true}\\[2pt]
\emptyout       & p(v)=\text{false}
\end{cases}
\]
$|\outfun|\in\{0,1\}$: cardinality can only shrink.

\subsection*{\texttt{take} : $\mathbb{N}^{+}\to T^{*}\to T^{*}$ (bounded by $n$)}
State $\mathrm{Counting}(k)$, $0\le k<n$:
\[
\outfun(\mathrm{Counting}(k),\nx{v}) =
\begin{cases}
\outseq{\nx{v}}        & k+1<n\\[2pt]
\outseq{\nx{v},\cmpl}  & k+1=n
\end{cases}
\]
At the boundary $|\outfun|=2$: value and synthesized completion.

\subsection*{\texttt{scan} vs. \texttt{reduce}}
With accumulator $a\in A$ and $\mathrm{acc}:A\to T\to A$:
\[
\outfun_{\text{scan}}(a,\nx{v}) = \outseq{\nx{\mathrm{acc}(a,v)}},
\qquad
\outfun_{\text{reduce}}(a,\nx{v}) = \emptyout,\quad
\outfun_{\text{reduce}}(a,\cmpl)  = \outseq{\nx{a},\cmpl}.
\]
Hence $\texttt{scan}:T^{*}\to A^{*}$ but $\texttt{reduce}:T^{*}\to A!$ (collapse at completion).

\section{Higher-order operators : $(T\to U^{*})\to T^{*}\to U^{*}$}
Augment the alphabet with tagged inner events
$\{\mathtt{in}(w,i),\mathtt{ic}(i),\mathtt{ie}(i)\}$. For \texttt{mergeMap}, with state
$(I,\mathit{flag})$, $I\in\mathcal{P}(\mathit{InnerId})$:
\[
\outfun((I,\mathit{flag}),\nx{v})=\emptyout,\quad
\outfun((I,\mathit{flag}),\mathtt{in}(w,i))=\outseq{\nx{w}},\quad
\outfun((\emptyset,\mathrm{SrcDone}),\mathtt{ic}(i))=\outseq{\cmpl}.
\]
Completion is non-local: it requires $\mathrm{SrcDone}\wedge I=\emptyset$. \texttt{switchMap}
replaces the inner id on each $\nx{v}$ (cancellation); \texttt{concatMap} carries a FIFO queue;
\texttt{exhaustMap} ignores $\nx{v}$ while an inner is live.

\section{Time-based operators (clock-augmented)}
With $\tick$ a unit tick and duration $d$:
\[
\begin{aligned}
\texttt{debounceTime}(d):\ &
\outfun((p,1,s),\tick)=\outseq{\nx{p}} &&(T^{*}\to T^{*},\ T!\to T?)\\
\texttt{bufferTime}(d):\ &
\outfun((B,1,\mathrm{Open}),\tick)=\outseq{\nx{B}} &&(T^{*}\to (T[])^{*})\\
\texttt{windowTime}(d):\ &
\outfun((w,1,\mathrm{Open}),\tick)=\outseq{\mathtt{wc}(w),\nx{w'}} &&(T^{*}\to (T^{*})^{+})
\end{aligned}
\]
For \texttt{windowTime} the single tick emits two symbols, which is why the outer cardinality
grows.

\section{Correspondence theorem}
\begin{theorem}[Static projection of dynamic length, informal]
Let $L(\mathrm{op})=\{\,|\outfun(s,\nx{v})| : (s,\nx{v})\text{ reachable}\,\}$. Then:
\begin{center}
\begin{tabular}{ll}
\toprule
Condition & C$\omega$ result \\
\midrule
$L=\{1\}$, $\cmpl\mapsto\outseq{\cmpl}$ & preserved ($T^{*}\to U^{*}$) \\
$L\subseteq\{0,1\}$                     & shrinks ($T!\to T?$) \\
$\sup L\ge 2$ or inner-driven           & grows (flattening) \\
$L=\{0\}$, $\cmpl\mapsto\outseq{\nx{a},\cmpl}$ & collapses to $A!$ \\
\bottomrule
\end{tabular}
\end{center}
\end{theorem}

\end{document}
